\section{Further research questions and a proposed program}
\label{sec:research-program}

The results suggest a research program centered on geometric interfaces and
explicit parameter budgets. The priorities below are ordered by their ability
to connect the proved local reductions to a complete recognition procedure.
Each asks for a mathematical statement or a reproducible artifact that can be
accepted or falsified independently.

\subsection{Priority I: close the native geometric interfaces}

\paragraph{1. A faithful through-tile producer.}
Can the source-parallelity charts and ordered cut-boundary words produce a
flat descriptor of the actual transported cubic pattern in polynomial encoded
time? The new theorem allows disconnected tiles and initially different
component partitions. A first target is a nontrivial complete class of normal
cuts with a proof that every tile component meets both interfaces. The producer
must record the spherical embedding, actual port identities, edge ownership,
and why subsequent context has no access to private vertices. A useful
experiment logs these facts and each precise rejection reason on native
hierarchy cuts; a count of large multiplicities alone tests none of them.

\paragraph{2. Compressed relative bands with marked transport.}
Can the proper-boundary handle theorem be made constructive at the level of
source charts without writing one band per unit of genus? Seek a compressed
relative spine whose expansion gives the actual band system and whose actions
on ordered boundary arcs are polynomial in the encoded input. One should
specify which operations on that spine remain polynomial: intersection,
concatenation, extraction of a co-core disc, and transport of a boundary
pattern. A scalar genus counter settles the unmarked census but is insufficient
for any operation that must return an embedded curve or disc.

\paragraph{3. Compile the saturated residue geometrically.}
Each saturated bordered base contributes one two-handle in a suitable relative
decomposition. Can its attaching curve be encoded directly from the ordered
interface words and bundle charts, with a short embedded certificate? A
natural first family is planar saturated bases with a bounded number of
interfaces, followed by orientable bases with large binary genus and then
nonorientable bases. The output must preserve the attaching map and framing;
recording only its abelianization loses the distinction exhibited by
$\langle x,y\mid x^3=y^2\rangle$. A successful theorem should bound both
the number of residual curves and their encoded interaction with the core.

\paragraph{4. Extract a source-addressed geometric noose.}
The graph kernel gives at most three original crossing edges and a compressed
shore coloring. Can the actual spherical embedding turn this into a simple
closed curve with a compressed isotopy certificate, then transport the
corresponding disc through the hierarchy? This is an appropriate use of
existing efficient algorithms for curves on surfaces
\cite{lackenby-curves}. The extra issue is the source provenance linking those
surface calculations to the three-dimensional cut and its successive boundary
identifications. The first deliverable should cover a repeated native product
region, including reversed charts and corners.

\subsection{Priority II: prove discovery and execution budgets}

\paragraph{5. Bound the cancellation component count on useful sources.}
The fifteen-tetrahedron example proves that a one-component assumption cannot
be imposed on every vertex meridian. Does some broader small bound hold after
a certified source transformation, or can one construct a family in which the
minimum possible component count tends to infinity? A convincing negative
family must control \emph{all} vertex meridians, rather than exhibit one
complicated meridian. A positive theorem should bound the expression
$r\log n+k\log(d+1)$, which is what the discovery algorithm actually pays.
This is a more flexible target than independently minimizing all three
parameters.

\paragraph{6. Control incidence without losing the one-vertex model.}
Can certified moves reduce the maximum tetrahedron-incidence degree $d$ while
keeping $t$ and the required cancellation budget under control? Subdivision
can lower local incidence in some representations but may introduce vertices;
the fixed-cocycle proof then stops applying. The desired statement must
quantify this tradeoff and preserve the marked exterior. A useful intermediate
theorem would bound the sum of local profile costs over all affected
tetrahedron stars after a Pachner move, rather than restart every profile
search on the whole source.

\paragraph{7. Several vertices with shared potentials.}
For multiple vertices, write a representative as $c+\delta\phi$ and account
for the sign chambers of its edge values. Can local cancellation profiles be
coupled through a small exact system in the shared vertex potentials? A
Schur-complement or arrangement-based approach is plausible when the number
of independent potentials is small. The proof must retain integer feasibility
and the correct boundary-cohomology quotient. A real-valued local solution or
an arbitrary fixed vector $|c|$ would not be a complete substitute. The earlier
sign-chamber machinery provides a baseline against which a new bound should
be compared \cite{incoming-cancellation}.

\paragraph{8. Reduce the active-system dependence on incidence.}
The factor $(Cd^2)^k$ counts connected supports and then local active
hyperplane choices. Could these two enumerations be combined through an
incidence graph of equations and edges? A potential gain is a smaller
dependence on $d$ or a bound in the number of distinct local coefficient
patterns. The difficulty is that an active-system incidence graph can
disconnect even when the cancellation support is connected. Any pruning rule
must retain all extremal profiles, including the homogeneous exceptional
case; the current complete algorithm is an exact oracle for small tests.

\paragraph{9. Independently certify restricted exhaustion.}
The positive rank gate is compact. A general negative search transcript is
not yet supplied. Can one certify that every connected support and every
relevant active system up to the chosen budgets has been accounted for,
sharing repeated coefficient patterns and exact row reductions? The one-edge
inventory for the retained counterexample is a tractable starting case.
Such a transcript would establish exhaustion of a \emph{restricted} disc
class. It would become a nontriviality certificate only after a separate
theorem makes that class sufficient for all required meridians.

\paragraph{10. An amortized hierarchy potential for all work.}
Which geometric potential bounds both useful cuts and corrective operations,
while a branch budget controls all alternatives actually explored? A candidate
must be tested against source rebuilding, pattern normalization, saturated
repairs, and failed discovery. The execution-tree statement in
Lemma~\ref{lem:branch-budget} provides an explicit acceptance criterion.
Instrumenting these charges is useful before proving a global bound: a local
strict descent that triggers expensive unrestricted restarts may have little
practical value despite its mathematical validity.

\subsection{Priority III: enlarge and improve the graph kernels}

\paragraph{11. Tiles with births, deaths, and private components.}
The through hypothesis excludes a component that meets only one interface or
neither. Is there a polynomial marked-cut kernel for every binary tile power
at a fixed threshold? One possible route uses the larger partition-diagram
semigroup and records components that become internal during composition.
Purely private repeated connected components need only a bounded number of
representatives to supply the two marks, but components that enter and later
leave a through region require a source-addressed transient analysis. A proof
must track those components; the rank recurrence of
Lemma~\ref{lem:through-stabilization} deliberately does not cover them.

\paragraph{12. Can the quadratic reblocking bound be reduced?}
The exponent bound $m\leq d$ and the universal connected cutoff are sharp.
They do not show that a $d^2Q$ kernel is necessary. The sharper estimate
$(a_k+1)\sum_{\mathcal O}h_{\mathcal O}Q_{\mathcal O}+(m-1)Q$ already
distinguishes easy instances. Can shared cycle structure give an $O(Q)$ or
$O(Q\log Q)$ explicit marked representative, or is a superlinear lower bound
forced by exact boundary marks and source lifting? Known terminal mimicking
networks suggest comparison tools \cite{chalermsook2021,liu2023}, but their
construction from an expanded graph does not by itself solve the compressed
input problem.

\paragraph{13. Stable kernels under nested descriptions.}
Flat powers have explicit tiles. A general straight-line grammar may duplicate
already compressed subgraphs, so repeated local kernelization can increase
intermediate size. Identify a grammar class closed under the required
geometric operations for which kernel size is bounded by a fixed polynomial
in the whole grammar. A useful first result would charge each original edge
definition only a logarithmic number of times. The proof must preserve exposed
vertex identity and exact mark unions at each composition, not only the final
yes/no answer.

\paragraph{14. Smaller positive graph proofs.}
The dense replay path intentionally makes a conservative quartic pass over
the finite kernel. A sparse replay can remove an avoidable factor, and
positive anchor-pair proofs may share cut or flow structure. The mathematical
question is whether a short family of certified witnesses implies coverage
of all anchor pairs without enumerating them individually. The practical
question is where that proof work dominates source construction on native
patterns. Changes should be measured on complete producer-plus-replay queries,
including small instances where reduction overhead loses.

\subsection{A concrete sequence of research deliverables}

\begin{enumerate}
\item \textbf{Native interface census.} Instrument a fixed corpus of actual
hierarchy cuts. Record $(t,k,r,d)$ where measurable, saturated versus proper
base attachments, faithful tile sizes, through failures, and all source and
replay costs. Retain inconclusive outcomes and exact source identities.
\item \textbf{One complete geometric adapter.} Choose the largest clearly
specified family in that census for which a faithful pattern grammar and
relative band transport can both be proved. Deliver a source-bound producer
and independent replay, then measure completed recognition stages.
\item \textbf{One global parameter theorem or obstruction family.} Establish
a sufficient discovery budget on a certified family, or construct a family
forcing its violation. Either result would materially refine the strategy;
more favorable timings without such structural information would not.
\item \textbf{A complete execution account.} Integrate the proved adapters
with the maintained recognizer, charge every failed attempt and repair, and
state the final complexity theorem with all input-size conversions. Until
that account is complete, retain the local results as proved reductions and
their global use as an explicit research objective.
\end{enumerate}
